% !TEX root = ../jlt.tex
\section{Guarantees from System~F}
\label{sec:guarantees}
% The three theorems about the translation of Section~\ref{sec:translation}
% and the appendices, stated after the exhibits so that each can
% point at one.  Proofs in Appendix~\ref{app:fjlt:theorems}; the
% adequacy argument (C++ -> FJlt) is Section~\ref{sec:translation:adequacy}
% and is not a theorem.

The theorems rest on FJlt alone.
Their proofs (Appendix~\ref{app:fjlt:theorems}) are inductions over the grammar, the typing and the translation of FJlt, and never mention the three rules.
The rules are not the hypotheses of the theorems.
They are what makes a \Cpp{} program a term of the calculus the theorems are about, and that step is the argument of Section~\ref{sec:translation:adequacy}.

\begin{theorem}[Coherence]
\label{thm:coherence}
The translation is a function of the program: two resolutions of the same constraint at the same concrete types yield the same record.
\end{theorem}

In \Cpp{} this is overload resolution together with the one-definition rule: at concrete types a specialisation is unique, or the program is ill-formed.
It gives ``the unit of \cppinline{+} at \cppinline{size_t}'' one meaning in a program, so a law can be stated about it, and it lets laws compose: two overloads written apart are chosen the same way wherever they are called, as the two behind Listing~\ref{lst:laws} are.

\begin{theorem}[Type preservation]
\label{thm:preservation}
If $\Delta; \Gamma \vdash e : \tau$ in FJlt, then $\llbracket \Delta \rrbracket; \llbracket \Gamma \rrbracket \vdash \llbracket e \rrbracket : \llbracket \tau \rrbracket$ in the target, System~F with records and iterators (Appendix~\ref{app:fjlt:translation}).
\end{theorem}

Each \cppinline{requires} clause has become a record parameter.
At each call, the instance the compiler resolved supplies a record of the type that parameter demands.

\begin{theorem}[Simulation]
\label{thm:simulation}
If $e \longrightarrow e'$ in FJlt (Appendix~\ref{app:fjlt:evaluation}), then $\llbracket e \rrbracket \longrightarrow^{+} \llbracket e' \rrbracket$ in the target, up to the reduction of a fold's count, a descending self-call's fuel and a node's reading (Appendix~\ref{app:fjlt:theorems}), and a value of base type translates to a constant.
For closed terms of base type, $e$ evaluates to $v$ if and only if $\llbracket e \rrbracket$ reduces to $\llbracket v \rrbracket$.
\end{theorem}

What a \cppinline{static_assert} sees is what the translation computes, and, by Rule~1, what the program computes at run time.
Section~\ref{sec:binding} is this theorem run twice.
The theorem also makes Section~\ref{sec:posture}'s equality exact: two closed terms of a carrier are equal when they evaluate to the same value, a constant of a base type compared bit for bit, and Leibniz's law of Rule~1 holds by substitution.

\begin{corollary}[Termination]
\label{cor:termination}
Every FJlt program terminates.
\end{corollary}

It follows from Theorem~\ref{thm:simulation} and strong normalisation of the target (Appendix~\ref{app:fjlt:translation}).
Every step of the program is at least one step of its image, and the image cannot step forever.
The count in Rule~3 is therefore not a guard against non-termination.
It is how a bound is said.

\begin{corollary}[Parametricity]
\label{cor:parametricity}
A generic FJlt function satisfies the free theorem of its System~F type.
\end{corollary}

By Theorem~\ref{thm:preservation} and Reynolds' abstraction theorem~\cite{reynolds1983types,wadler1989theorems}.
\cppinline{concat} at a monoid can only combine the words by \texttt{op}, starting from \texttt{unit}, which is the free theorem of a fold.
It is also why an instance with a context (Section~\ref{sec:translation:contexts}) can read a composite's law off its components' records.
The composite cannot do anything the components' records do not provide.

\begin{listing}[htbp]
\begin{minipage}[t]{0.49\linewidth}
\begin{cpp}
// the overload associativity licenses
template <typename T, typename Op,
          IsCountWord N>
  requires IsMonoid<T, Op>
constexpr T power(const T& x, N n, Op op) {
  if (n == 0) {
    return unit<T>(op);
  } else {  // the halves are equal: one, twice
    const T h = power(x, n / 2, op);
    const T hh = op(h, h);
    return n % 2 == 1 ? op(hh, x) : hh;
  }
}
template <IsBoundedRange R,
          typename T = ranges::range_value_t<R>,
          typename Op = plus<T>>
  requires IsMonoid<T, Op> && IsCopies<R>
constexpr T concat(R r, Op op = {}) {
  const auto n = copies_count(r);
  return n == 0 ? unit<T>(op)
                : power(copies_value(r), n, op);
}
\end{cpp}
\end{minipage}\hfill
\begin{minipage}[t]{0.49\linewidth}
\begin{cpp}
// the action law, declared
using Act = remove_cvref_t<decltype(act)>;
template <> struct acting_monoid<Act> {
  using type = multiplies<M>;
};
// the overload it licenses
template <IsBoundedRange R,
          semiregular M, typename H>
  requires IsCopies<R> && IsRightAction<
               H, M, ranges::range_value_t<R>>
constexpr M fold(R r, M v, H h) {
  using Op = acting_monoid<H>::type;
  return h(v, concat(r, Op{}));
}
\end{cpp}
\end{minipage}
\begin{cpp}
// Listing 3's fold, by the library: the laws choose how
constexpr size_t fibonacci(size_t n) {
  return fold(n >> chain<size_t>() >> constant<size_t>(fib), V{1, 0}, act)[1];
}
static_assert(fibonacci(10) == fs_10[1] &&
              fibonacci(100'000'000) == 14'139'011'350'745'967'675uz);
\end{cpp}
\begin{minted}{text}
fibonacci = λ n. π₂ (fold act (1, 0) ms)      -- ms = map (const M) (chain prior n)
          = λ n. π₂ (act (1, 0) (concat (·) ms))      -- by the action law
          = λ n. π₂ (act (1, 0) (Mⁿ))                 -- by associativity, squaring
\end{minted}
\caption{Laws composing, on Listing~\ref{lst:rule3}'s chain.
Top left: the monoid concatenation $M \cdot M \cdots M$ is replaced by the power $M^n$, implemented by repeated squaring~\cite[\S 4.6.3]{knuth1997art2}.
Top right: the type constraints license the power method (Rule~2 for the specialisation, Theorem~\ref{thm:coherence} for the composition).
Middle: the power method keeps the $10^8$ terms of the Fibonacci sequence within the compile-time step limit (Table~\ref{tab:binding}).
Bottom: the System~F translation closely resembles the \Cpp{} above (Figure~\ref{fig:arrows}).}
\label{lst:laws}
\end{listing}

\begin{corollary}[Type soundness]
\label{cor:soundness}
A well-typed closed FJlt term of base type evaluates to a value.
\end{corollary}

This is the theorem a reader of Featherweight Java~\cite{igarashi2001featherweight} expects first, soundness, here stronger because every term ends, and cheap once the rest is in place.
A closed well-typed term is a value or takes a step (progress, Lemma~\ref{lem:progress}), and it cannot step forever (Corollary~\ref{cor:termination}), so it reaches a value.
What the translation adds is the second half of Theorem~\ref{thm:simulation}: the value reached is the value the System~F image computes.

\paragraph{Certified and inferred laws.}
A program declares few laws and infers the rest.
A declared law is an axiom, its certificate evaluated where a record is built.
A law at a construction is derived from the laws of its parts, through an instance with a context, by the entailment of qualified types~\cite{jones1994qualified}: that a finite power of a semiring is a semimodule over it is inferred, not declared (Appendix~\ref{app:fjlt:inferred}).
The inference is local in Pierce and Turner's sense~\cite{pierce2000local}: information moves only between adjacent nodes, and here only upward, from a term's immediate parts to the term, a synthesised attribute~\cite{knuth1968semantics}.
They infer type arguments and Jlt infers laws, so the likeness is in the flow of information, not in what is inferred.
By Theorem~\ref{thm:coherence} each inferred law is unique, and the conditions on contexts that end instance search~\cite{sulzmann2007understanding} end the inference.
As type checking rules out classes of errors without proving a program correct, law inference licenses classes of optimisations without deciding program equivalence, which stays undecidable where every program ends: a program testing whether a machine halts within $n$ steps equals the constant false exactly when that machine never halts~\cite{turing1936computable}.
Each optimisation is local, a rewrite at the node whose law licenses it, and so mechanical; the global one, the best equivalent program, is the undecidable one.
Since the laws are the application's, not the machine's, so are the optimisations they license.

